\documentclass[11pt]{amsart}
\pdfsuppressptexinfo=-1
\usepackage[margin=1.1in]{geometry}
\usepackage[T1]{fontenc}
\usepackage{lmodern}
\usepackage{amsmath,amssymb,amsthm,mathtools}
\usepackage{booktabs,enumitem,array}
\usepackage[section]{placeins}
\usepackage{xcolor}
\usepackage[hidelinks]{hyperref}
\hypersetup{pdftitle={A lower bound for the Bohr radius of the bidisc},
  pdfauthor={Tejasvi Singh Tomar}}

\newtheorem{theorem}{Theorem}[section]
\newtheorem{lemma}[theorem]{Lemma}
\newtheorem{proposition}[theorem]{Proposition}
\newtheorem{corollary}[theorem]{Corollary}
\theoremstyle{remark}
\newtheorem{remark}[theorem]{Remark}
\numberwithin{equation}{section}
\setlength{\emergencystretch}{2em}
\raggedbottom

\newcommand{\C}{\mathbb{C}}
\newcommand{\D}{\mathbb{D}}
\newcommand{\T}{\mathbb{T}}
\newcommand{\Pclass}{\mathcal{P}_2}
\newcommand{\Bfun}{\mathcal{B}}
% ---------------------------------------------------------------------------------------------
% Numerical values (rounded in the safe direction).
\newcommand{\Kbound}{0.301665}
\newcommand{\rstar}{\frac{60333}{200000}}
\newcommand{\ParA}{\frac{20674053}{62500000}}
\newcommand{\ParB}{\frac{3901457}{1562500000}}
\newcommand{\ParG}{\frac{101990031}{10^9}}
\newcommand{\SidonRoot}{0.301339}
\newcommand{\KneseVal}{0.3006}
\newcommand{\KneseRoot}{0.3006282\ldots}
\newcommand{\PatelUB}{0.302825279492}
\newcommand{\PatelLB}{0.30158758}
\newcommand{\BPWWUB}{0.3177}
\newcommand{\GriegoUB}{0.3174541}
\newcommand{\YeVal}{2.186529771892569\ldots}
\newcommand{\FiveOwn}{2.187}
\newcommand{\NBoxes}{8{,}912}
\newcommand{\NVert}{269{,}769}
\newcommand{\TailN}{40}
\newcommand{\Evalue}{0.99999556}
\newcommand{\SixOwn}{2.363}
\newcommand{\NBoxesSix}{61{,}118}
\newcommand{\NVertSix}{9{,}893{,}956}
\newcommand{\MinSlackSix}{$8.48\cdot10^{-8}$}
\newcommand{\EvaluePatel}{0.99999997857}
\newcommand{\MinSlack}{$1.49\cdot10^{-5}$}
\newcommand{\lean}[1]{\texttt{\detokenize{#1}}}
\newcommand{\LeanBound}{0.301657}
\newcommand{\SUPPDOI}{10.5281/zenodo.23716440}
\newcommand{\REPOURL}{https://github.com/tejasvi/bidisc-bohr-radius/tree/v6}
\newcommand{\SevenOwn}{2.552}
\newcommand{\SevenNum}{2.55166}
\newcommand{\NFilesSeven}{576}
\newcommand{\NBoxesSeven}{210{,}960}
\newcommand{\NVertSeven}{65{,}041{,}100}
\newcommand{\MinSlackSeven}{$8.88\cdot10^{-8}$}
\newcommand{\SizeSeven}{$26.8$}
% ---------------------------------------------------------------------------------------------

\begin{document}

\title[The Bohr radius of the bidisc]{A lower bound for the Bohr radius of the bidisc}
\author{Tejasvi Singh Tomar}
\renewcommand{\urladdrname}{{\itshape ORCID}}
\urladdr{\href{https://orcid.org/0000-0003-4668-3639}{https://orcid.org/0000-0003-4668-3639}}
\subjclass[2020]{32A05; 30B10, 42A05, 47A48, 68V05, 68V20}
\keywords{Bohr radius, bidisc, Carath\'eodory class, Sidon constant, transfer-function realization, computer-assisted proof, formal verification}

\begin{abstract}
We prove that the Bohr radius $K_2$ of the bidisc, whose exact value is unknown, satisfies $K_2>\Kbound$, improving Patel's bound $K_2>\PatelLB$. As in his argument, we pass to functions with positive real part and value $1$ at the origin; we bound the sum of the moduli of their coefficients of degree $n$ by $2S_n$ instead of $2\sqrt n$, where $S_n$ is the Sidon constant of $\{0,1,\dots,n\}$, and we bound $S_5$, $S_6$ and $S_7$ with computer-assisted certificates. In Lean~4 we also prove $K_2\ge\Kbound$, assuming Patel's inequality coupling the coefficients of degrees two and three.
\end{abstract}

\maketitle

\section{Introduction}

For $d\ge1$ let $\D^d=\{z\in\C^d:|z_1|,\dots,|z_d|<1\}$ be the unit polydisc and let the Schur class $\mathcal S(\D^d)$ be the set of holomorphic maps $f\colon\D^d\to\overline{\D}$. For $f(z)=\sum_\alpha c_\alpha z^\alpha$ and $r>0$ write $\|f_r\|_1=\sum_\alpha|c_\alpha|r^{|\alpha|}$, where $|\alpha|=\alpha_1+\dots+\alpha_d$. The Bohr radius of $\D^d$ is
\[
K_d=\sup\{r>0:\|f_r\|_1\le1\text{ for all }f\in\mathcal S(\D^d)\}.
\]
Bohr \cite{Boh14} proved that $K_1\ge1/6$. The sharp value $K_1=1/3$ was found independently by Riesz, Schur and Wiener, and Bohr's paper includes Wiener's proof (see \cite{BK97}). Boas and Khavinson \cite{BK97} introduced $K_d$ for $d\ge2$, proved $1/(3\sqrt d)<K_d<2\sqrt{(\log d)/d}$ and asked for the exact value of $K_d$ when $d>1$. Defant, Frerick, Ortega-Cerd\`a, Ouna\"ies and Seip \cite{DFOOS11} showed that $K_d$ has order $\sqrt{(\log d)/d}$, and Bayart, Pellegrino and Seoane-Sep\'ulveda \cite{BPS14} proved $K_d\sim\sqrt{(\log d)/d}$ as $d\to\infty$. No exact value of $K_d$ is known for any $d\ge2$.

For the bidisc, functions of one variable show $K_2\le1/3$. Knese \cite[Corollary~1.2]{Kne25} proved $K_2\ge\KneseVal$: by his Theorem~1.1 and And\^o's inequality, $K_2$ is at least the root $\KneseRoot$ of $\sum_{k\ge1}\sqrt k\,r^k=1/2$. Baran, Pikul, Woerdeman and Wojtylak \cite[Theorem~6.4]{BPWW26} proved $K_2<\BPWWUB$, and Griego \cite{Gri26} lowered this to $K_2<\GriegoUB$.

In an unrefereed solution to a problem on MathDB, Patel \cite{PatMDB} proved $K_2>\PatelLB$. He also reports the upper bound $K_2<\PatelUB$, from a rational inner function and an exact computation that his reproducibility package \cite{PatZen} carries out; we have not checked this bound, but if it holds, it and Theorem~\ref{thm:main} confine $K_2$ to an interval of width about $1.2\cdot10^{-3}$.

Sidon constants enter through the coefficients. For a finite set $\Lambda\subset\mathbb Z$ the Sidon constant $S(\Lambda)$ is the least $C$ with
\[
\sum_{\lambda\in\Lambda}|a_\lambda|\le C\max_{|w|=1}\Bigl|\sum_{\lambda\in\Lambda}a_\lambda w^\lambda\Bigr|\qquad(a_\lambda\in\C),
\]
and we write
\[
S_k=S(\{0,1,\dots,k\})\qquad(k\ge1).
\]
Bernal-Gonz\'alez, Cabana, Garc\'ia, Maestre, Mu\~noz-Fern\'andez and Seoane-Sep\'ulveda \cite[Theorem~3.3]{BCGMMS21} showed that $K_d$ is at least the root of $\sum_{k\ge1}S(k,d)\,r^k=1/2$, where $S(k,d)$ is the Sidon constant of the $k$-homogeneous polynomials in $d$ variables. For $d=2$, $S(k,2)=S_k$ (Remark~\ref{rem:skd}). Newman's bound $S_k\le\sqrt k$, whose proof appears in Shapiro's thesis \cite[Chapter~3]{Sha51} (see also \cite[(1)]{Neu26}), turns this equation into Knese's, and Shapiro \cite{Sha51} showed that $S_k=\sqrt k$ if and only if $k\in\{1,2,4\}$. Few values of $S_k$ are known; Table~\ref{tab:sidon} lists them with their sources.

We give Patel's argument for the lower bound in full, confirm the exact rational computation behind it with our own programs, and improve the bound.

\begin{theorem}\label{thm:main}
The Bohr radius of the bidisc satisfies $K_2>\Kbound$.
\end{theorem}

The gain over Patel's bound, about $8\cdot10^{-5}$, comes from the terms of degrees five, six and seven.

Patel reduces the problem to the coefficients of holomorphic functions $H$ on $\D^2$ with $\operatorname{Re}H>0$ and $H(0)=1$ (Section~\ref{sec:setup}). Carath\'eodory's inequality on slices bounds sums of moduli of the coefficients of degree two and, with the help of a positive semidefinite Toeplitz form, of degree three, and a third inequality (Lemma~\ref{lem:couple}) couples the two degrees; Patel combines the three into a bound for the coefficients of degrees two and three together (Section~\ref{sec:patel}). The third inequality comes from the unitary realization of Schur functions on the bidisc \cite{Agl90,BT98}, which rests on And\^o's inequality \cite{And63}. In the degrees $n\ge4$ Patel uses the bound $2\sqrt n$ for the sum of the absolute values of the coefficients of degree $n$ of $H$, which he derives from Knese's estimate \cite[Lemma~2.1]{Kne25}.

Our improvement starts from a simple observation: Carath\'eodory's inequality on slices also bounds this sum by $2S_n$ (Lemma~\ref{lem:slice}). By Newman's bound and Shapiro's theorem, $S_n<\sqrt n$ for every $n\ge5$, while $S_4=2$. We prove $S_5\le\FiveOwn$, $S_6\le\SixOwn$ and $S_7\le\SevenOwn$ with computer-assisted certificates built from positive definite Toeplitz matrices (Section~\ref{sec:values}); for $S_5$ our bound is slightly above the value $S_5=\YeVal$ that Ye et al.\ determined \cite{Ye26}. In the degrees $n\ge8$ we keep Newman's bound $S_n\le\sqrt n$, since the size of the certificates grows quickly with $n$ (Table~\ref{tab:certs}). With rational upper bounds for $\sqrt n$ and an explicit tail, an exact rational computation shows that Patel's majorant, with these bounds, is less than $1$ at $r=\Kbound$, and by continuity also at slightly larger $r$, which proves Theorem~\ref{thm:main} (Section~\ref{sec:proof}). Section~\ref{sec:verif} describes the computer checks and a formal proof in Lean, conditional on Lemma~\ref{lem:couple}, of the non-strict bound $K_2\ge\Kbound$, which needs only the computation at $r=\Kbound$.

\section{The Carath\'eodory normalization and slices}\label{sec:setup}

Let the Carath\'eodory class $\Pclass$ be the set of holomorphic functions $H$ on $\D^2$ with $\operatorname{Re}H>0$ and $H(0)=1$. We write
\[
H(z)=1+\sum_{n\ge1}H_n(z),\qquad H_n(z)=\sum_{j+k=n}h_{jk}z_1^jz_2^k,
\]
and set $\|H_n\|_1=\sum_{j+k=n}|h_{jk}|$ and $M_r(H)=\sum_{n\ge1}\|H_n\|_1\,r^n$.

We use the following reduction.

\begin{lemma}\label{lem:norm}
Let $0<r<1$. If $M_r(H)\le1$ for every $H\in\Pclass$, then $K_2\ge r$.
\end{lemma}

\begin{proof}
Constant functions satisfy $\|f_r\|_1\le1$. Let $f\in\mathcal S(\D^2)$ be nonconstant. By the maximum principle $|f|<1$ on $\D^2$. Multiplying $f$ by a unimodular constant does not change $\|f_r\|_1$, so we may assume that $t=f(0)$ lies in $[0,1)$. Then $H=(1-f)/(1-t)$ lies in $\Pclass$. Its coefficients of positive degree are those of $f$ multiplied by $-1/(1-t)$. Hence $\|f_r\|_1=t+(1-t)M_r(H)\le1$.
\end{proof}

\begin{remark}\label{rem:converse}
The normalization loses nothing: if $r<K_2$, then $M_r(H)\le1$ for every $H\in\Pclass$. Indeed, for $0<s<1$ the function $z\mapsto H(sz)$ maps $\overline{\D}^2$ into a compact subset of the right half-plane, hence into a disc $\{w:|w-\rho|<\rho\}$ with $\rho\ge1$. Then $f(z)=1-H(sz)/\rho$ lies in $\mathcal S(\D^2)$ with $f(0)=1-1/\rho$, and $\|f_r\|_1\le1$ becomes $\sum_n\|H_n\|_1\,s^nr^n\le1$. Let $s\to1$. Thus $K_2=\sup\{r:M_r(H)\le1\text{ for all }H\in\Pclass\}$.
\end{remark}

\begin{lemma}\label{lem:slice}
Let $H\in\Pclass$ and $n\ge1$. Then $|H_n(\zeta)|\le2$ for $\zeta\in\T^2$, and $\|H_n\|_1\le2S_n$.
\end{lemma}

\begin{proof}
Fix $\zeta\in\T^2$. The function $\varphi(\lambda)=H(\lambda\zeta)=1+\sum_{m\ge1}H_m(\zeta)\lambda^m$ is holomorphic on $\D$ with $\operatorname{Re}\varphi>0$. For $0<\rho<1$,
\[
H_n(\zeta)\rho^n=\frac1\pi\int_0^{2\pi}\operatorname{Re}\varphi(\rho e^{i\theta})\,e^{-in\theta}\,d\theta,
\]
so $|H_n(\zeta)|\rho^n\le2\operatorname{Re}\varphi(0)=2$. Letting $\rho\to1$, we obtain Carath\'eodory's inequality $|H_n(\zeta)|\le2$. Now write $H_n(\zeta_1,\zeta_2)=\zeta_2^n\,q_n(\zeta_1/\zeta_2)$ with $q_n(w)=\sum_{j=0}^nh_{j,n-j}w^j$. As $\zeta$ runs over $\T^2$, the quotient $\zeta_1/\zeta_2$ runs over $\T$ and $|H_n(\zeta)|=|q_n(\zeta_1/\zeta_2)|$. Hence $\max_\T|q_n|\le2$, and by the definition of $S_n$,
\[
\|H_n\|_1=\sum_{j=0}^n|h_{j,n-j}|\le S_n\max_\T|q_n|\le2S_n.\qedhere
\]
\end{proof}

\begin{remark}\label{rem:skd}
The same identification shows that $S(k,2)=S_k$: on $\T^2$ a $k$-homogeneous polynomial in two variables is, up to the factor $\zeta_2^k$, a polynomial in $\zeta_1/\zeta_2$ of degree at most $k$. With this identification, the next corollary is the case $d=2$ of \cite[Theorem~3.3]{BCGMMS21}.
\end{remark}

\begin{corollary}\label{cor:sidon}
If $0<r<1$ and $\sum_{k\ge1}S_kr^k\le1/2$, then $K_2\ge r$.
\end{corollary}

\begin{proof}
By Lemma~\ref{lem:slice}, $M_r(H)\le2\sum_{k\ge1}S_kr^k\le1$ for every $H\in\Pclass$. Apply Lemma~\ref{lem:norm}.
\end{proof}

In \cite{BCGMMS21} the corresponding estimate is proved for $f$ itself. There Wiener's inequality takes the place of Carath\'eodory's: the Taylor coefficients $\hat g(k)$ of a holomorphic $g\colon\D\to\overline{\D}$ satisfy $|\hat g(k)|\le1-|g(0)|^2$ for $k\ge1$. Patel derives the bound $\|H_n\|_1\le2\sqrt n$ for all $H\in\Pclass$ from Knese's estimate \cite[Lemma~2.1]{Kne25}, by a limiting argument with $(1-\varepsilon H)/(1+\varepsilon H)$ \cite{PatMDB}. Lemma~\ref{lem:slice} with Newman's bound gives it directly.

With the values of $S_k$ for $k\le5$ in Table~\ref{tab:sidon} and $S_k\le\sqrt k$ for $k\ge6$, Corollary~\ref{cor:sidon} gives $K_2\ge\SidonRoot$. If the value of Ye et al.\ for $S_5$ and Newman's bounds for $S_6$ and $S_7$ are replaced by the bounds of Proposition~\ref{prop:s5}, it gives $K_2\ge0.301370$. Patel's coupling of degrees two and three goes beyond this.

\section{The coefficients of degrees two and three}\label{sec:patel}

Throughout this section $H\in\Pclass$, and
\[
A_2=|h_{20}|+|h_{02}|,\quad B_2=|h_{11}|,\qquad A_3=|h_{30}|+|h_{03}|,\quad B_3=|h_{21}|+|h_{12}|,
\]
so that $\|H_2\|_1=A_2+B_2$ and $\|H_3\|_1=A_3+B_3$. Inner products are linear in the first argument, $\langle x,y\rangle=y^*x$.

\begin{lemma}[Patel {\cite{PatMDB}}]\label{lem:quad}
$A_2^2+B_2^2\le4$.
\end{lemma}

\begin{proof}
Choose $\zeta_0\in\T$ such that $h_{20}\zeta_0^2$ and $h_{02}$ have the same argument, and put $a=h_{20}\zeta_0^2+h_{02}$ and $b=h_{11}\zeta_0$. Then $|a|=A_2$, $|b|=B_2$ and $H_2(\pm\zeta_0,1)=a\pm b$. By Lemma~\ref{lem:slice}, $|a\pm b|\le2$, and therefore $|a|^2+|b|^2=\tfrac12(|a+b|^2+|a-b|^2)\le4$.
\end{proof}

\begin{lemma}\label{lem:toeplitz}
Let $R$ be a positive semidefinite Hermitian Toeplitz matrix of size $k+1$, and let $q(\zeta)=\sum_{j=0}^kc_j\zeta^j$. Then $c^*Rc\le R_{00}\max_\T|q|^2$, where $c=(c_0,\dots,c_k)^{\mathsf T}$.
\end{lemma}

\begin{proof}
Let $e_0,\dots,e_k$ be vectors in $\C^{k+1}$ with $\langle e_l,e_j\rangle=R_{jl}$, for example the columns of $R^{1/2}$. Since $R$ is Toeplitz, the tuples $(e_0,\dots,e_{k-1})$ and $(e_1,\dots,e_k)$ have the same Gram matrix. Hence $e_j\mapsto e_{j+1}$ $(0\le j<k)$ extends to a well-defined linear isometry from the span of $e_0,\dots,e_{k-1}$ onto the span of $e_1,\dots,e_k$, and then to a unitary operator $\Omega$ on the span of all the $e_j$. Then $e_j=\Omega^je_0$, so $\sum_jc_je_j=q(\Omega)e_0$. As $\Omega$ is unitary, $\|q(\Omega)\|\le\max_\T|q|$, and
\[
c^*Rc=\Bigl\|\sum_jc_je_j\Bigr\|^2=\|q(\Omega)e_0\|^2\le\max_\T|q|^2\,\|e_0\|^2=R_{00}\max_\T|q|^2.\qedhere
\]
\end{proof}

Lemma~\ref{lem:toeplitz} is a consequence of the Carath\'eodory--Toeplitz theorem. Ye et al.\ prove it in the same way \cite[Lemma~3.1]{Ye26}.

\begin{lemma}[Patel {\cite{PatMDB}}]\label{lem:cubic}
Let $p(\zeta)=\sum_{j=0}^3c_j\zeta^j$ satisfy $\max_\T|p|\le1$. Then
\[
35\bigl(|c_0|+|c_3|\bigr)^2+28\bigl(|c_1|+|c_2|\bigr)^2\le44.
\]
Consequently $35A_3^2+28B_3^2\le176$.
\end{lemma}

\begin{proof}
Replacing $p(\zeta)$ by $p(e^{i\theta}\zeta)$ for a suitable $\theta$, we may assume that $c_1$ and $c_2$ have the same argument. Choose $\xi\in\T$ with $|\bar\xi c_0+c_3|=|c_0|+|c_3|$, and let
\[
u=(3,4,-4,-3)^{\mathsf T},\qquad v=(0,1,1,0)^{\mathsf T},\qquad w=(\xi,0,0,1)^{\mathsf T},\qquad R=uu^*+28vv^*+35ww^*.
\]
The matrix $R$ is positive semidefinite. Since $|\xi|=1$, it is also Toeplitz, with first row $(44,12,-12,-9+35\xi)$. By Lemma~\ref{lem:toeplitz}, with $c=(c_0,\dots,c_3)^{\mathsf T}$,
\[
|u^*c|^2+28|v^*c|^2+35|w^*c|^2=c^*Rc\le44.
\]
Here $|v^*c|=|c_1+c_2|=|c_1|+|c_2|$ and $|w^*c|=|\bar\xi c_0+c_3|=|c_0|+|c_3|$, which proves the first claim. For the second, apply it to $p=q_3/2$ from the proof of Lemma~\ref{lem:slice}. Then $|c_0|+|c_3|=A_3/2$ and $|c_1|+|c_2|=B_3/2$.
\end{proof}

The proof above extends to a one-parameter family. For $-1\le t\le2$ and $u_t=(t,1,-1,-t)^{\mathsf T}$, the matrix $u_tu_t^*+(t+1)vv^*+(2+t-t^2)ww^*$ is positive semidefinite and Toeplitz, with first row $(2+t,t,-t,-t^2+(2+t-t^2)\xi)$, so the same argument shows
\[
(2+t-t^2)\bigl(|c_0|+|c_3|\bigr)^2+(t+1)\bigl(|c_1|+|c_2|\bigr)^2\le2+t.
\]
Lemma~\ref{lem:cubic} is the case $t=3/4$. The case $t=1/2$ is the inequality $9(|c_0|+|c_3|)^2+6(|c_1|+|c_2|)^2\le10$, from which Yu et al.\ obtain $S_3\le5/3$ by the Cauchy--Schwarz inequality \cite[(6)]{Yu26}. On its own, Lemma~\ref{lem:cubic} gives only $\|H_3\|_1\le\sqrt{396/35}\approx3.364$, more than $2S_3=10/3$. It is used together with the next inequality, which bounds $A_3$ in terms of $B_2$.

\begin{lemma}[Patel {\cite{PatMDB}}]\label{lem:couple}
$A_3+\tfrac12B_2^2\le2$.
\end{lemma}

\begin{proof}
We use the transfer-function realization of Schur functions on the bidisc. The Schur--Agler class consists of the $g\in\mathcal S(\D^2)$ with $\|g(X_1,X_2)\|\le1$ for every pair of commuting strict contractions $X_1$, $X_2$ on a Hilbert space. By And\^o's inequality \cite{And63}, every $g\in\mathcal S(\D^2)$ belongs to this class. By Agler's theorem \cite{Agl90} (see also \cite{BT98,AM02}) there are a Hilbert space $\mathcal K$, complementary orthogonal projections $P$ and $Q$ on $\mathcal K$, and an isometry $V=\begin{psmallmatrix}\kappa&x^*\\y&W\end{psmallmatrix}$ of $\C\oplus\mathcal K$, the colligation, such that
\[
g(z)=\kappa+x^*Z\bigl(I-WZ\bigr)^{-1}y,\qquad Z=z_1P+z_2Q\qquad(z\in\D^2).
\]
For $z\in\D^2$ we have $\|Z\|\le\max_j|z_j|<1$, so $I-WZ$ and the other operators $I-ZW$, $I-ZU$ inverted below are invertible.
We may take $V$ unitary, which matters if the realization is given with an isometric colligation, as in \cite[Section~2]{Kne25}. Indeed, the operator $\begin{psmallmatrix}V&I-VV^*\\0&V^*\end{psmallmatrix}$ on $(\C\oplus\mathcal K)\oplus(\C\oplus\mathcal K)$ is unitary and agrees with $V$ on the first summand. We use it as the new colligation, with state space $\mathcal K\oplus(\C\oplus\mathcal K)$ and projections $P\oplus I$ and $Q\oplus0$. Since the new $y$ is the old one, the new $W$ and $x^*$ agree with the old ones on $\mathcal K$, and the new projections map $\mathcal K$ into itself, each term $x^*Z(WZ)^my$ of the expansion of $g$ is unchanged. From now on $\mathcal K$, $V$, $x$, $y$, $W$, $P$ and $Q$ denote the new data.

Apply this to $g=(H-1)/(H+1)$, which lies in $\mathcal S(\D^2)$ and vanishes at $0$. Then $\kappa=0$, and since $V$ is unitary,
\[
\|x\|=\|y\|=1,\qquad Wx=0,\qquad W^*y=0,\qquad xx^*+W^*W=I,\qquad yy^*+WW^*=I.
\]
Put $U=W+yx^*$. These identities imply $U^*U=UU^*=I$. Let $G=I-ZW$. Then $Gx=x$ and $x^*G^{-1}Zy=x^*Z(I-WZ)^{-1}y=g$. Since $I-ZU=G-Zyx^*$ is a rank-one perturbation of $G$ and $1-g\ne0$,
\[
x^*(I-ZU)^{-1}x=x^*G^{-1}x+\frac{(x^*G^{-1}Zy)(x^*G^{-1}x)}{1-x^*G^{-1}Zy}=1+\frac{g}{1-g}=\frac1{1-g}.
\]
Hence $x^*(I+ZU)(I-ZU)^{-1}x=2x^*(I-ZU)^{-1}x-1=(1+g)/(1-g)=H$, and expanding the resolvent,
\[
H_n(z)=2x^*(ZU)^nx\qquad(n\ge1).
\]
Relative to $\mathcal K=P\mathcal K\oplus Q\mathcal K$ write
\[
U=\begin{pmatrix}U_{11}&U_{12}\\U_{21}&U_{22}\end{pmatrix},\qquad x=\begin{pmatrix}a\\b\end{pmatrix},\qquad Ux=\begin{pmatrix}\tilde a\\\tilde b\end{pmatrix},
\]
so that $\|a\|^2+\|b\|^2=\|\tilde a\|^2+\|\tilde b\|^2=1$. Collecting the monomials $z_1z_2$, $z_1^3$ and $z_2^3$ in $2x^*(ZU)^nx$, we find
\[
\tfrac12h_{11}=a^*U_{12}\tilde b+b^*U_{21}\tilde a,\qquad\tfrac12h_{30}=a^*U_{11}^2\tilde a,\qquad\tfrac12h_{03}=b^*U_{22}^2\tilde b.
\]
Put $\sigma=\|U_{11}^*a\|^2+\|U_{22}^*b\|^2$ and $\tau=\|U_{11}\tilde a\|^2+\|U_{22}\tilde b\|^2$. Since $U^*$ is an isometry, $\|U_{11}^*a\|^2+\|U_{12}^*a\|^2=\|a\|^2$ and $\|U_{21}^*b\|^2+\|U_{22}^*b\|^2=\|b\|^2$, so $\|U_{12}^*a\|^2+\|U_{21}^*b\|^2=1-\sigma$. Since $U$ is an isometry, $\|U_{12}\tilde b\|^2+\|U_{21}\tilde a\|^2=1-\tau$ in the same way. Applying the Cauchy--Schwarz inequality to $\tfrac12h_{11}=\langle\tilde b,U_{12}^*a\rangle+\langle\tilde a,U_{21}^*b\rangle=\langle U_{12}\tilde b,a\rangle+\langle U_{21}\tilde a,b\rangle$ in the two groupings, we obtain
\[
\tfrac14B_2^2\le1-\sigma\qquad\text{and}\qquad\tfrac14B_2^2\le1-\tau.
\]
Finally $\tfrac12|h_{30}|=|\langle U_{11}\tilde a,U_{11}^*a\rangle|\le\tfrac12(\|U_{11}^*a\|^2+\|U_{11}\tilde a\|^2)$, and likewise for $h_{03}$, so $A_3\le\sigma+\tau\le2-\tfrac12B_2^2$.
\end{proof}

\begin{remark}
Lemmas~\ref{lem:quad} and~\ref{lem:couple} are sharp. Equality holds in both for $H=(1+z_1)/(1-z_1)$, where $A_2=A_3=2$ and $B_2=0$, and for $H=(1+z_1z_2)/(1-z_1z_2)$, where $A_2=A_3=0$ and $B_2=2$.
\end{remark}

\begin{lemma}[Patel {\cite{PatMDB}}]\label{lem:combine}
Let $\alpha,\beta>0$, $\gamma\ge0$ and $r>0$, and put
\[
\Bfun(r)=4\alpha+176\beta+2\gamma+\frac1{4\alpha}+\frac1{4\alpha+2\gamma}+\frac{(r-\gamma)^2}{140\beta}+\frac{r^2}{112\beta}.
\]
Then $\|H_2\|_1+r\|H_3\|_1\le\Bfun(r)$.
\end{lemma}

\begin{proof}
Expanding the squares verifies the identity
\begin{align*}
\Bfun(r)-\bigl(A_2+B_2+r(A_3+B_3)\bigr)={}&\alpha\bigl(4-A_2^2-B_2^2\bigr)+\beta\bigl(176-35A_3^2-28B_3^2\bigr)+\gamma\bigl(2-A_3-\tfrac12B_2^2\bigr)\\
&+\alpha\Bigl(A_2-\frac1{2\alpha}\Bigr)^2+\Bigl(\alpha+\frac\gamma2\Bigr)\Bigl(B_2-\frac1{2\alpha+\gamma}\Bigr)^2\\
&+35\beta\Bigl(A_3-\frac{r-\gamma}{70\beta}\Bigr)^2+28\beta\Bigl(B_3-\frac r{56\beta}\Bigr)^2.
\end{align*}
By Lemmas~\ref{lem:quad}, \ref{lem:cubic} and~\ref{lem:couple}, each term on the right is nonnegative.
\end{proof}

Thus $\Bfun(r)$ is the Lagrangian bound, with multipliers $\alpha$, $\beta$ and $\gamma$, for the maximum of $A_2+B_2+r(A_3+B_3)$ under the three inequalities of Lemmas~\ref{lem:quad}, \ref{lem:cubic} and~\ref{lem:couple}.

\section{Sidon constants and certificates for \texorpdfstring{$S_5$, $S_6$ and $S_7$}{small k}}\label{sec:values}

The proof of Theorem~\ref{thm:main} uses upper bounds for $S_1$ and for $S_k$ with $k\ge4$. Degrees two and three are covered by Lemma~\ref{lem:combine}.

\begin{table}[ht]
\centering
\caption{Sidon constants $S_k=S(\{0,1,\dots,k\})$.}\label{tab:sidon}
\begin{tabular}{@{}lll@{}}
\toprule
$k$ & $S_k$ & source\\
\midrule
$1$ & $1$ & elementary\\
$2$ & $\sqrt2$ & Newman (see \cite{Neu01})\\
$3$ & $5/3$ & Yu et al.\ \cite{Yu26}; lower bound Neuwirth \cite{Neu26}\\
$4$ & $2$ & Newman's bound; equality by Shapiro \cite{Sha51}\\
$5$ & $\YeVal$; $\le\FiveOwn$ & Ye et al.\ \cite{Ye26}; Proposition~\ref{prop:s5}\\
$6$ & $\le\SixOwn$ & Proposition~\ref{prop:s5}\\
$7$ & $\le\SevenOwn$ & Proposition~\ref{prop:s5}\\
$k\ge8$ & $\le\sqrt k$ & Newman \cite[Chapter~3]{Sha51}\\
\bottomrule
\end{tabular}
\end{table}

The proof of Theorem~\ref{thm:main} uses neither the value $S_3=5/3$ nor the value of $S_5$ below, so that it does not rely on these recent preprints. Ye et al.\ \cite[Theorem~1.1]{Ye26} show that $S_5$ is the root in $(2.18,2.19)$ of $675X^4-1026X^3-936X^2-414X+677$. The proof of Proposition~\ref{prop:s5} rests on the following lemma.

\begin{lemma}\label{lem:dual}
Let $k\ge1$. For $w\in\C^{k+1}$ put
\[
\nu(w)=\max\Bigl\{\operatorname{Re}\sum_{j=0}^k\bar w_jc_j:\ \max_\T\Bigl|\sum_{j=0}^kc_jz^j\Bigr|\le1\Bigr\}.
\]
\begin{enumerate}[label=\textup{(\alph*)}]
\item $S_k=\max\{\nu(w):w_0=w_1=1,\ |w_j|=1\ (2\le j\le k)\}$.
\item $\nu$ is convex.
\item Let $y\in\C^k$, and let $T(y)$ be the Hermitian Toeplitz matrix with diagonal entries $1$ and entries $T(y)_{jl}=y_{l-j}$ for $l>j$. If $T(y)$ is positive definite, then $\nu(w)^2\le w^*T(y)^{-1}w$.
\item For $0<\delta\le\pi/4$, the arc $\{e^{i\varphi}:|\varphi-\varphi_0|\le\delta\}$ lies in the triangle with vertices $e^{i(\varphi_0-\delta)}$, $e^{i(\varphi_0+\delta)}$ and $e^{i\varphi_0}/\cos\delta$.
\item $\nu(\bar w)=\nu(w)$, where $\bar w=(\bar w_0,\dots,\bar w_k)$.
\end{enumerate}
\end{lemma}

\begin{proof}
(a) Write $q(z)=\sum_jc_jz^j$. If $|w_j|=1$ for all $j$, then $\nu(w)\le\max\{\sum_j|c_j|:\max_\T|q|\le1\}=S_k$. Equality holds when $c$ attains this maximum and $w_j=c_j/|c_j|$ (any unimodular number if $c_j=0$). The substitution $c_j\mapsto c_je^{i(\psi+j\theta)}$ replaces $q(z)$ by $e^{i\psi}q(e^{i\theta}z)$ and $w_j$ by $w_je^{-i(\psi+j\theta)}$ without changing $\nu$, and a choice of $\psi$ and $\theta$ makes $w_0=w_1=1$.
(b) $\nu$ is a supremum of real-linear functions of $w$.
(c) By Lemma~\ref{lem:toeplitz}, $c^*T(y)c\le\max_\T|q|^2$ for $c=(c_0,\dots,c_k)^{\mathsf T}$. Hence, by the Cauchy--Schwarz inequality,
\[
\operatorname{Re}w^*c\le\bigl(w^*T(y)^{-1}w\bigr)^{1/2}\bigl(c^*T(y)c\bigr)^{1/2}\le\bigl(w^*T(y)^{-1}w\bigr)^{1/2}\max_\T|q|.
\]
(d) After a rotation, $\varphi_0=0$. The triangle is $\{u\in\C:\operatorname{Re}u\ge\cos\delta,\ \operatorname{Re}(ue^{-i\delta})\le1,\ \operatorname{Re}(ue^{i\delta})\le1\}$, and $u=e^{i\varphi}$ with $|\varphi|\le\delta$ satisfies $\cos\varphi\ge\cos\delta$ and $\cos(\varphi\mp\delta)\le1$.
(e) The map $c\mapsto\bar c$ maps the set of admissible $c$ onto itself, since $\sum_j\bar c_jz^j=\overline{q(\bar z)}$, and $\operatorname{Re}\sum_jw_j\bar c_j=\operatorname{Re}\sum_j\bar w_jc_j$. Hence $\nu(\bar w)=\max\operatorname{Re}\sum_jw_jc_j=\max\operatorname{Re}\sum_jw_j\bar c_j=\nu(w)$.
\end{proof}

Parts (a) and (c) are variants of the reduction to phase functionals and of the Toeplitz domination bound used by Ye et al.\ \cite[(14) and Proposition~3.2]{Ye26}, who normalize the phases differently \cite[(17)]{Ye26}. Proposition~\ref{prop:s5} also differs from their proof in the covering of the phase space.

\begin{proposition}\label{prop:s5}
$S_5\le\FiveOwn$, $S_6\le\SixOwn$ and $S_7\le\SevenOwn$.
\end{proposition}

\begin{proof}[Proof \textup{(computer-assisted)}]
Let $(k,b_k)$ be $(5,\FiveOwn)$, $(6,\SixOwn)$ or $(7,\SevenOwn)$. By Lemma~\ref{lem:dual}(a) it suffices to show that $\nu(1,1,e^{i\varphi_2},\dots,e^{i\varphi_k})\le b_k$ for $(\varphi_2,\dots,\varphi_k)\in[-\pi,\pi]^{k-1}$. By Lemma~\ref{lem:dual}(e), which replaces $(\varphi_2,\dots,\varphi_k)$ by $(-\varphi_2,\dots,-\varphi_k)$, we may also assume $\varphi_2\ge0$. For $k=5$ and $k=6$ the certificate covers $[-\pi,\pi]^{k-1}$ by boxes, obtained from a uniform grid of $8^{k-1}$ boxes with half-widths $\pi/8$ by repeated bisection. For $k=7$ it covers $[0,\pi]\times[-\pi,\pi]^5$ in the same way, starting from a uniform grid of $2\cdot4^5$ boxes with half-widths $\pi/4$. By Lemma~\ref{lem:dual}(d), the vectors $(1,1,e^{i\varphi_2},\dots,e^{i\varphi_k})$ with $\varphi$ in the box lie in a product of $k-1$ triangles. By Lemma~\ref{lem:dual}(b), the maximum of $\nu$ over this product is attained at one of its $3^{k-1}$ vertices. For each vertex $v$ that occurs, the certificate gives $y\in\C^k$ with $T(y)$ positive definite and $v^*T(y)^{-1}v<b_k^2$, so that $\nu(v)<b_k$ by Lemma~\ref{lem:dual}(c). A vertex with a coordinate $-1$ occurs with two labels, phase $-\pi$ and phase $\pi$, and has a certificate for each. Table~\ref{tab:certs} gives the sizes of the three certificates, and Section~\ref{sec:verif} describes the check.
\end{proof}

\begin{table}[ht]
\centering
\caption{The certificates of Proposition~\ref{prop:s5}.}\label{tab:certs}
\begin{tabular}{@{}llll@{}}
\toprule
 & $k=5$ & $k=6$ & $k=7$\\
\midrule
bound $b_k$ & $\FiveOwn$ & $\SixOwn$ & $\SevenOwn$\\
initial grid & $8^4$ & $8^5$ & $2\cdot4^5$\\
boxes & $\NBoxes$ & $\NBoxesSix$ & $\NBoxesSeven$\\
distinct vertex labels & $\NVert$ & $4{,}333{,}585$ & $46{,}344{,}804$\\
smallest margin $b_k^2-v^*T(y)^{-1}v$ & \MinSlack & \MinSlackSix & \MinSlackSeven\\
\bottomrule
\end{tabular}
\end{table}

Numerical maximization, which is not part of the proof, gives $S_6\approx2.36256$ and $S_7\approx\SevenNum$, so the bounds $\SixOwn$ and $\SevenOwn$ are close to the true values. Further gains along these lines need bounds for $S_n$ with $n\ge8$ or sharper inequalities in the degrees two and three.

\section{Proof of Theorem~\ref{thm:main}}\label{sec:proof}

Throughout the paper, lower bounds are rounded down and upper bounds up. Let $N=\TailN$. Let $s_4=2$, $s_5=\FiveOwn$, $s_6=\SixOwn$, $s_7=\SevenOwn$, and $s_n=\lceil10^8\sqrt n\,\rceil/10^8$ for $8\le n\le N$. By Table~\ref{tab:sidon} and Proposition~\ref{prop:s5}, $S_n\le s_n$ for $4\le n\le N$, and $S_n\le\sqrt n\le n$ for all $n$. Let $H\in\Pclass$, let $0<r<1$, and fix $\alpha,\beta>0$ and $\gamma\ge0$. Lemma~\ref{lem:slice} bounds $\|H_n\|_1$ by $2$ for $n=1$, by $2s_n$ for $4\le n\le N$ and by $2n$ for $n>N$, and Lemma~\ref{lem:combine} bounds $\|H_2\|_1+r\|H_3\|_1$. Summing the series $\sum_{n>N}nr^n=r^{N+1}(N+1-Nr)/(1-r)^2$, we obtain the majorant
\[
M_r(H)\le E(r):=2r+r^2\Bfun(r)+2\sum_{n=4}^{N}s_nr^n+\frac{2r^{N+1}(N+1-Nr)}{(1-r)^2}.
\]
We take
\[
r_*=\rstar,\qquad\alpha=\ParA,\qquad\beta=\ParB,\qquad\gamma=\ParG.
\]
All quantities in $E(r_*)$ are then rational, and an exact computation shows $E(r_*)<\Evalue<1$. Since $\alpha$, $\beta$ and $\gamma$ are fixed, $E$ is continuous on $(0,1)$, so there is $r'>r_*$ with $E(r')<1$. By Lemma~\ref{lem:norm}, $K_2\ge r'>r_*$, which proves Theorem~\ref{thm:main}.

\begin{remark}[Patel's computation]\label{rem:patel}
With Patel's parameters $\alpha=264631/800000$, $\beta=249627/10^8$ and $\gamma=10197021/10^8$, with $N=20$ and with $s_n=\lceil10^8\sqrt n\,\rceil/10^8$ for $4\le n\le20$, the function $E$ is the rational majorant that he evaluates in \cite{PatMDB}; its value at his radius $r_P=15079379/(5\cdot10^7)$ is less than $\EvaluePatel$, below his bound $1-1/(5\cdot10^7)$.
\end{remark}

\section{Verification}\label{sec:verif}

\subsection{Computer checks}\label{sec:compchecks}
The certificates of Proposition~\ref{prop:s5} are checked in the ball arithmetic of Arb \cite{Joh17}, through python-flint \cite{pyflint}. The checker reads the numbers of the certificates as binary64 numbers, exactly, rebuilds the boxes from the recorded bisections and verifies at every vertex the condition stated in the proof of Proposition~\ref{prop:s5}; for $S_6$ and $S_7$, whose certificates are split into files, a further program checks, from the outputs of the file checks, that the files together cover the phase domain. The majorant $E(r_*)$ of Section~\ref{sec:proof} is evaluated in exact rational arithmetic, and so is Patel's majorant in Remark~\ref{rem:patel}. The programs and their output are archived with the certificates (see Data availability).

\subsection{Formal proof}\label{sec:lean}

The formal proof is written in the Lean~4 proof assistant \cite{Lean4} with the Mathlib library \cite{Mathlib}. The formal statements use Mathlib's definitions. A family $c\colon\mathbb N^2\to\C$ is admissible if the series $\sum_\alpha c_\alpha z^\alpha$ converges absolutely at every point of $\D^2$ and its value there has modulus less than $1$. Bohr's inequality holds at $r$ if $\sum_{\alpha\in F}|c_\alpha|r^{|\alpha|}\le1$ for every admissible family $c$ and every finite set $F\subset\mathbb N^2$, and $K_2$ is the supremum of the radii $r>0$ at which it holds. It agrees with the definition in the introduction: a nonconstant $f\in\mathcal S(\D^2)$ maps $\D^2$ into $\D$, and a constant of modulus $1$ satisfies Bohr's inequality with equality. The main theorem is the following.
\begin{theorem}[formalized in Lean]\label{thm:formal}
Suppose that $A_3+\tfrac12B_2^2\le2$, the inequality of Lemma~\ref{lem:couple}, holds for every family $h$ with $h_{00}=1$ whose series converges absolutely on $\D^2$ to values with positive real part, that is, for the coefficients of every $H\in\Pclass$. Then $K_2\ge\Kbound$.
\end{theorem}
The proof of the hypothesis rests on And\^o's inequality and Agler's realization, which are not in Mathlib. The formal proof follows Sections~\ref{sec:setup} to~\ref{sec:proof} and checks the certificates of Proposition~\ref{prop:s5} in one of two ways.
\begin{enumerate}[label=(\arabic*),leftmargin=*]
\item A checking program written in Lean validates the certificates, converted into Lean source files, in exact rational arithmetic and ball arithmetic, and a Lean proof shows that every part it accepts implies the bound of Lemma~\ref{lem:dual}(c) on the boxes of that part, so that the converted files need not be trusted. These checks run as compiled code, through \lean{native_decide}, so the trusted base includes Lean's compiler and runtime; the rest of the proof uses only the standard axioms.
\item A second formalization proves the three bounds of Proposition~\ref{prop:s5} with every check run by Lean's kernel. With the step of the first proof that derives the main theorem from them, the main theorem follows from four theorems in which every check is run by the kernel; these are not combined into a single Lean theorem.
\end{enumerate}

\section*{Data availability}
The certificates of Proposition~\ref{prop:s5}, the programs of Section~\ref{sec:verif}, the programs that produced the certificates, the controls and the output of all runs, and the Lean files of Section~\ref{sec:lean}, are archived at \href{https://doi.org/\SUPPDOI}{doi:\SUPPDOI}. The same files, except the $48$ files of the certificate for $S_6$, the $\NFilesSeven$ files of the certificate for $S_7$ and an archive of generated Lean files of the kernel-only proof of $S_7\le\SevenOwn$, are also in the repository \url{\REPOURL}. The README of the archive lists the Lean theorems and projects, with the versions of Lean and Mathlib, and gives the commands that run the checks and build the proofs. The certificate for $S_5$ has the SHA-256 checksum
\begin{center}
\texttt{ce3e9dd41a0dbd9c3a8ea49b04aa7aaf21f0fe557291a2084c7739461a460a14},
\end{center}
and the archive lists the checksums of the other files.

\section*{Use of AI tools}
Generative AI tools were used to find and prove the results, to write the programs, and to draft the text. The author reviewed the arguments and the programs and takes responsibility for the content.

\begin{thebibliography}{99}

\bibitem{Agl90} J.~Agler, \emph{On the representation of certain holomorphic functions defined on a polydisc}, in: Topics in Operator Theory: Ernst D.~Hellinger Memorial Volume, Oper. Theory Adv. Appl. \textbf{48}, Birkh\"auser, Basel, 1990, 47--66.

\bibitem{AM02} J.~Agler and J.~E.~McCarthy, \emph{Pick interpolation and Hilbert function spaces}, Grad. Stud. Math. \textbf{44}, Amer. Math. Soc., Providence, RI, 2002.

\bibitem{And63} T.~And\^o, \emph{On a pair of commutative contractions}, Acta Sci. Math. (Szeged) \textbf{24} (1963), 88--90.

\bibitem{BT98} J.~A.~Ball and T.~T.~Trent, \emph{Unitary colligations, reproducing kernel Hilbert spaces, and Nevanlinna--Pick interpolation in several variables}, J. Funct. Anal. \textbf{157} (1998), 1--61.

\bibitem{BPWW26} R.~Baran, P.~Pikul, H.~J.~Woerdeman and M.~Wojtylak, \emph{Contractive realization theory for the annulus and other intersections of disks on the Riemann sphere}, J. Funct. Anal. \textbf{290} (2026), no.~8, 111346; arXiv:2504.03236.

\bibitem{BPS14} F.~Bayart, D.~Pellegrino and J.~B.~Seoane-Sep\'ulveda, \emph{The Bohr radius of the $n$-dimensional polydisk is equivalent to $\sqrt{(\log n)/n}$}, Adv. Math. \textbf{264} (2014), 726--746.

\bibitem{BCGMMS21} L.~Bernal-Gonz\'alez, H.~J.~Cabana, D.~Garc\'ia, M.~Maestre, G.~A.~Mu\~noz-Fern\'andez and J.~B.~Seoane-Sep\'ulveda, \emph{A new approach towards estimating the $n$-dimensional Bohr radius}, Rev. R. Acad. Cienc. Exactas F\'is. Nat. Ser. A Mat. RACSAM \textbf{115} (2021), no.~2, Paper No.~44.

\bibitem{BK97} H.~P.~Boas and D.~Khavinson, \emph{Bohr's power series theorem in several variables}, Proc. Amer. Math. Soc. \textbf{125} (1997), 2975--2979.

\bibitem{Boh14} H.~Bohr, \emph{A theorem concerning power series}, Proc. London Math. Soc. (2) \textbf{13} (1914), 1--5.

\bibitem{Car24} M.~Carneiro, \emph{Lean4Lean: Verifying a typechecker for Lean, in Lean}, arXiv:2403.14064, 2024.

\bibitem{DFOOS11} A.~Defant, L.~Frerick, J.~Ortega-Cerd\`a, M.~Ouna\"ies and K.~Seip, \emph{The Bohnenblust--Hille inequality for homogeneous polynomials is hypercontractive}, Ann. of Math. (2) \textbf{174} (2011), 485--497.

\bibitem{Gri26} S.~Griego, \emph{Improve $C_{59}$ upper bound}, pull request no.~75 to the repository teorth/optimizationproblems, GitHub, May 2026, \url{https://github.com/teorth/optimizationproblems/pull/75}.

\bibitem{Joh17} F.~Johansson, \emph{Arb: efficient arbitrary-precision midpoint-radius interval arithmetic}, IEEE Trans. Comput. \textbf{66} (2017), 1281--1292.

\bibitem{Kne25} G.~Knese, \emph{Three radii associated to Schur functions on the polydisk}, Proc. Amer. Math. Soc. Ser. B \textbf{12} (2025), 48--63; arXiv:2410.21693v3.

\bibitem{Lean4} L.~de~Moura and S.~Ullrich, \emph{The Lean 4 theorem prover and programming language}, in: Automated Deduction -- CADE 28, Lecture Notes in Comput. Sci. \textbf{12699}, Springer, 2021, 625--635.

\bibitem{Mathlib} The mathlib Community, \emph{The Lean mathematical library}, in: Proc. 9th ACM SIGPLAN Int. Conf. Certified Programs and Proofs (CPP 2020), 2020, 367--381.

\bibitem{nanoda} \lean{nanoda_lib}, a type checker for Lean~4, \url{https://github.com/ammkrn/nanoda_lib}.

\bibitem{Neu01} S.~Neuwirth, \emph{The Sidon constant of sets with three elements}, arXiv:math/0102145, 2001.

\bibitem{Neu26} S.~Neuwirth, \emph{On the (Fourier analytic) Sidon constant of $\{0,1,2,3\}$}, arXiv:2603.28229, 2026 (note dated 2013).

\bibitem{PatMDB} S.~Patel, \emph{Improved lower and upper bounds for the bidisc Bohr radius}, solution to the problem ``Exact value of the Bohr radius of the bidisc'', MathDB, second solution on the page, posted 27 August 2026 (UTC) and edited the same day, \url{https://mathdb.com/p/392062/exact-value-of-the-bohr-radius-of-the-bidisc#s2} (consulted 21 October 2026).

\bibitem{PatZen} S.~Patel, \emph{Complex phases and a certified upper bound for the bidisc Bohr radius: reproducibility package}, version 1.0.0, Zenodo, 5 September 2026, doi:10.5281/zenodo.22341928.

\bibitem{pyflint} python-flint: Python bindings for FLINT, version~0.9.0 (FLINT~3.6.0), \url{https://github.com/flintlib/python-flint}.

\bibitem{Sha51} H.~S.~Shapiro, \emph{Extremal problems for polynomials and power series}, Master's thesis, Massachusetts Institute of Technology, 1951.

\bibitem{Ye26} Z.~Ye et al., \emph{The exact Sidon constant of six consecutive frequencies}, preprint, October 2026, \url{https://github.com/ABCDEFGCoder/The-Sidon-Constant-of-n-Consecutive-Frequencies}.

\bibitem{Yu26} H.~Yu et al., \emph{The exact Sidon constant of $\{0,1,2,3\}$ and its formalization}, preprint, ai.viXra.org:2606.0009, 2026 (dated 10 May 2026); Lean formalization at \url{https://github.com/yhx-12243/Sidon3}.

\end{thebibliography}
\end{document}
